End-to-end autonomous driving is entering a new phase where learning-based systems directly map rich sensor streams to driving actions. Recent progress in end-to-end driving stacks—especially those powered by vision–language–action (VLA) models—suggests a promising path toward unified perception, prediction, and planning within a single scalable framework~\cite{intelligence2025pi05visionlanguageactionmodelopenworld,black2026pi0visionlanguageactionflowmodel,hwang2024emma}. These models benefit from large-scale data and strong representation learning, and they increasingly demonstrate competence in complex urban interactions. However, as end-to-end systems become more capable, a central bottleneck is no longer only model capacity, but the lack of a scalable and reliable evaluation and training infrastructure that matches the pace of progress.

A scalable real-world simulator is particularly critical in the end-to-end era. Unlike modular pipelines, where intermediate outputs can be unit-tested (e.g., detection AP, tracking MOTA, forecasting ADE~\cite{fong2021panoptic,nuscenes2019,WaymoWorldModel2026}), end-to-end policies are typically evaluated by closed-loop driving outcomes that are sensitive to subtle distribution shifts and compounding errors. Unfortunately, current evaluations for end-to-end systems still primarily rely on real-world testing, which is costly to run and difficult to scale. More importantly, real-world testing is often biased and under-covered—limited by geography, weather, traffic density, and the rarity of safety-critical events—and is therefore hard to reproduce and hard to compare fairly across methods. This evaluation gap slows iteration, obscures failure modes, and makes it challenging to establish trustworthy progress for end-to-end autonomy.

Beyond evaluation, an interactive real-world simulator is also a key enabler for effective online reinforcement learning (RL) for VLA systems. Online RL requires large volumes of closed-loop interactions and diverse counterfactual experiences—precisely the regimes where real-world driving is unsafe, expensive, and ethically constrained. A simulator that can faithfully model future sensor observations under proposed actions can unlock scalable training signals: policies can explore alternative maneuvers, learn to recover from off-nominal or risky states, and improve robustness under rare events, all while remaining in a controlled and repeatable environment. In this sense, simulation is not merely a testbed—it becomes the engine for continual improvement.

However, building a simulator suitable for end-to-end autonomy is challenging. It must generate photorealistic sensor observations while maintaining strict action-following (causal consistency with ego motion), cross-view consistency (agreement among multiple cameras), and long-horizon stability (avoiding drift over extended rollouts). Moreover, it should support controllability over dynamic agents and static road elements when optional conditions are available, enabling targeted scenario synthesis and stress-testing. Existing approaches often fall short along one or more of these axes: they may produce visually plausible videos but ignore fine-grained action causality, generate per-view videos that disagree geometrically, or degrade quickly when rolled out long-term. Many high-quality video diffusion models are also designed for offline, bidirectional generation over a full clip, which limits their applicability for real-time interaction.

In this work, we present \textbf{\modelname}, a controllable multi-view generative world model for autonomous driving built upon video diffusion. Given history multi-camera video streams and a driving action (or action sequence) to be applied, our model generates the corresponding future multi-camera video streams. When available, the model additionally accepts optional conditions to control dynamic traffic agents and static road elements, enabling targeted counterfactual rollouts and scenario editing. 
Besides the action and scene-element control modules, we retain a text-conditioning branch to control global appearance (\eg, weather and time of day), enabling controllable appearance-driven style editing of the generated videos.
Crucially, unlike traditional bidirectional video diffusion models, \modelname~operates in a streaming, autoregressive fashion, generating futures incrementally for real-time interaction. This design makes the model naturally compatible with closed-loop usage—both for scalable evaluation of end-to-end policies and for online RL training where the simulator must respond promptly to newly sampled actions.

Our approach achieves high-quality multi-view video generation with (i) high view-consistency, ensuring coherent geometry and object identity across cameras; (ii) strict action following, producing futures aligned with commanded ego behavior; and (iii) long video rollout ability, enabling stable predictions over extended horizons. Together, these properties move generative world models closer to a practical “real-world simulator” abstraction—one that can support reproducible benchmarking, scalable regression testing, and interactive learning for end-to-end/VLA driving systems.